\section{The alternating {\v C}ech complex}
\label{section-alternating-cech}

\noindent
Let $S$ be a scheme. Let $f : U \to X$ be an \'etale morphism of algebraic
spaces over $S$. The functor
$$
j : U_{spaces, \etale} \longrightarrow X_{spaces, \etale},\quad
V/U \longmapsto V/X
$$
induces an equivalence of $U_{spaces, \etale}$ with the localization
$X_{spaces, \etale}/U$, see
Properties of Spaces, Section \ref{spaces-properties-section-localize}.
Hence there exist functors
$$
f_! : \textit{Ab}(U_\etale) \longrightarrow
\textit{Ab}(X_\etale),\quad
f_! : \textit{Mod}(\mathcal{O}_U) \longrightarrow \textit{Mod}(\mathcal{O}_X),
$$
which are left adjoint to
$$
f^{-1} : \textit{Ab}(X_\etale) \longrightarrow
\textit{Ab}(U_\etale),\quad
f^* : \textit{Mod}(\mathcal{O}_X) \longrightarrow \textit{Mod}(\mathcal{O}_U)
$$
see
Modules on Sites, Section \ref{sites-modules-section-localize}.
Warning: This functor, a priori, has
nothing to do with cohomology with compact supports!
We dubbed this functor ``extension by zero'' in the reference above.
Note that the two versions of $f_!$ agree as $f^* = f^{-1}$ for
sheaves of $\mathcal{O}_X$-modules.

\medskip\noindent
As we are going to use this construction below let us recall some of its
properties. Given an abelian sheaf $\mathcal{G}$ on $U_\etale$
the sheaf $f_!$ is the sheafification of the presheaf
$$
V/X \longmapsto
f_!\mathcal{G}(V) =
\bigoplus\nolimits_{\varphi \in \Mor_X(V, U)}
\mathcal{G}(V \xrightarrow{\varphi} U),
$$
see
Modules on Sites, Lemma \ref{sites-modules-lemma-extension-by-zero}.
Moreover, if $\mathcal{G}$ is an $\mathcal{O}_U$-module, then $f_!\mathcal{G}$
is the sheafification of the exact same presheaf of abelian groups which
is endowed with an $\mathcal{O}_X$-module structure in an obvious way
(see loc.\ cit.). Let $\overline{x} : \Spec(k) \to X$
be a geometric point. Then there is a canonical identification
$$
(f_!\mathcal{G})_{\overline{x}} =
\bigoplus\nolimits_{\overline{u}} \mathcal{G}_{\overline{u}}
$$
where the sum is over all $\overline{u} : \Spec(k) \to U$ such that
$f \circ \overline{u} = \overline{x}$, see
Modules on Sites, Lemma \ref{sites-modules-lemma-stalk-j-shriek}
and
Properties of Spaces, Lemma
\ref{spaces-properties-lemma-points-small-etale-site}.
In the following we are going to study the sheaf $f_!\underline{\mathbf{Z}}$.
Here $\underline{\mathbf{Z}}$ denotes the constant sheaf on
$X_\etale$ or $U_\etale$.

\begin{lemma}
\label{lemma-product-is-tensor-product}
Let $S$ be a scheme. Let $f_i : U_i \to X$ be \'etale morphisms
of algebraic spaces over $S$. Then there are isomorphisms
$$
f_{1, !}\underline{\mathbf{Z}} \otimes_{\mathbf{Z}}
f_{2, !}\underline{\mathbf{Z}}
\longrightarrow
f_{12, !}\underline{\mathbf{Z}}
$$
where $f_{12} : U_1 \times_X U_2 \to X$ is the structure morphism
and
$$
(f_1 \amalg f_2)_! \underline{\mathbf{Z}}
\longrightarrow
f_{1, !}\underline{\mathbf{Z}} \oplus
f_{2, !}\underline{\mathbf{Z}}
$$
\end{lemma}

\begin{proof}
Once we have defined the map it will be an isomorphism by our description
of stalks above. To define the map it suffices to work on the level of
presheaves. Thus we have to define a map
$$
\left(\bigoplus\nolimits_{\varphi_1 \in \Mor_X(V, U_1)} \mathbf{Z}\right)
\otimes_{\mathbf{Z}}
\left(\bigoplus\nolimits_{\varphi_2 \in \Mor_X(V, U_2)} \mathbf{Z}\right)
\longrightarrow
\bigoplus\nolimits_{\varphi \in \Mor_X(V, U_1 \times_X U_2)}
\mathbf{Z}
$$
We map the element $1_{\varphi_1} \otimes 1_{\varphi_2}$ to the element
$1_{\varphi_1 \times \varphi_2}$ with obvious notation. We omit the proof
of the second equality.
\end{proof}

\noindent
Another important feature is the trace map
$$
\text{Tr}_f : f_!\underline{\mathbf{Z}} \longrightarrow \underline{\mathbf{Z}}.
$$
The trace map is adjoint to the
map $\mathbf{Z} \to f^{-1}\underline{\mathbf{Z}}$ (which is an isomorphism).
If $\overline{x}$ is above, then $\text{Tr}_f$ on stalks at $\overline{x}$
is the map
$$
(\text{Tr}_f)_{\overline{x}} :
(f_!\underline{\mathbf{Z}})_{\overline{x}} =
\bigoplus\nolimits_{\overline{u}} \mathbf{Z}
\longrightarrow
\mathbf{Z} = \underline{\mathbf{Z}}_{\overline{x}}
$$
which sums the given integers. This is true because it is adjoint to the map
$1 : \mathbf{Z} \to f^{-1}\underline{\mathbf{Z}}$. In particular, if
$f$ is surjective as well as \'etale then $\text{Tr}_f$ is surjective.

\medskip\noindent
Assume that $f : U \to X$ is a surjective \'etale
morphism of algebraic spaces. Consider the {\it Koszul complex}
associated to the trace map we discussed above
$$
\ldots \to \wedge^3f_!\underline{\mathbf{Z}} \to
\wedge^2f_!\underline{\mathbf{Z}} \to f_!\underline{\mathbf{Z}} \to
\underline{\mathbf{Z}} \to 0
$$
Here the exterior powers are over the sheaf of rings $\underline{\mathbf{Z}}$.
The maps are defined by the rule
$$
e_1 \wedge \ldots \wedge e_n \longmapsto
\sum\nolimits_{i = 1, \ldots, n} (-1)^{i + 1}
\text{Tr}_f(e_i)
e_1 \wedge \ldots \wedge \widehat{e_i} \wedge \ldots \wedge e_n
$$
where $e_1, \ldots, e_n$ are local sections of $f_!\underline{\mathbf{Z}}$.
Let $\overline{x}$ be a geometric point of $X$ and set
$M_{\overline{x}} = (f_!\underline{\mathbf{Z}})_{\overline{x}} =
\bigoplus_{\overline{u}} \mathbf{Z}$. Then the stalk of the complex above at
$\overline{x}$ is the complex
$$
\ldots \to \wedge^3 M_{\overline{x}} \to \wedge^2 M_{\overline{x}}
\to M_{\overline{x}} \to \mathbf{Z} \to 0
$$
which is exact because $M_{\overline{x}} \to \mathbf{Z}$ is surjective, see
More on Algebra, Lemma \ref{more-algebra-lemma-homotopy-koszul-abstract}.
Hence if we let $K^\bullet = K^\bullet(f)$ be the complex with
$K^i = \wedge^{i + 1}f_!\underline{\mathbf{Z}}$, then we obtain a
quasi-isomorphism
\begin{equation}
\label{equation-quasi-isomorphism}
K^\bullet \longrightarrow \underline{\mathbf{Z}}[0]
\end{equation}
We use the complex $K^\bullet$ to define what we call
the alternating {\v C}ech complex associated to $f : U \to X$.

